% Ad Atticum 15.3 (ad-atticum-15-03)
% Source: https://raw.githubusercontent.com/PerseusDL/canonical-latinLit/master/data/phi0474/phi057/phi0474.phi057.perseus-lat1.xml
% Fetched by scripts/fetch_latin.py

% Dateline (Perseus): Scr. in Arpinati xt K. Iun. a. 710 (44) .

\ciceroLetterOpener{CICERO ATTICO salutem}

\ciceroSection{1} undecimo Kal. accepi in Arpinati duas epistulas tuas, quibus duabus meis respondisti. una erat xv Kal., altera xii data. ad superiorem igitur prius. accurres in Tusculanum, ut scribis; quo me vi Kal. venturum arbitrabar. quod scribis parendum victoribus, non mihi quidem cui sunt multa potiora. nam illa quae recordaris Lentulo et Marcello consulibus acta in aede Apollinis, nec causa eadem est nec simile tempus, praesertim cum Marcellum scribas aliosque discedere. erit igitur nobis coram odorandum et constituendum tutone Romae esse possimus. novi conventus habitatores sane movent; in magnis enim versamur angustiis. sed sunt ista parvi; quin et maiora contemnimus. Calvae testamentum cognovi, hominis turpis ac sordidi; tabula Demonici quod tibi curae est gratum. de †malo† scripsi iam pridem ad Dolabellam accuratissime, modo redditae litterae sint. eius causa et cupio et debeo.

\ciceroSection{2} venio ad propiorem. cognovi de Alexione quae desiderabam. Hirtius est tuus. Antonio, †quoniam† est, volo peius esse. de Quinto filio, ut scribis, †A. M. C.† de patre coram agemus. Brutum omni re qua possum cupio iuvare. cuius de oratiuncula idem te quod me sentire video. sed parum intellego quid me velis scribere quasi a Bruto habita oratione, cum ille ediderit. qui tandem convenit? an sic ut in tyrannum iure optimo caesum? multa dicentur, multa scribentur a nobis sed alio modo et tempore. de sella Caesaris bene tribuni; praeclaros etiam xiv ordines! Brutum apud me fuisse gaudeo, modo et libenter fuerit et sat diu.
